% \documentclass[../main.tex]{subfiles}

% \begin{document}

\chapter{Supersymmetry}

\section{Quantum Field Theory and the Standard Model \textcolor{gray}{(Short: enough QFT to discuss Feynman gauge/lack of Goldstone bosons and enough SM to motivate SUSY as extension and introduce relevant definitions from electroweak theory)}}

\section{Constructing a supersymmetric theory}
\subsection{Superspace}
\subsection{Supersymmetric Lagrangians}

\subsubsection{The supersymmetric Lagrangian \textcolor{red}{FROM SUSY PROJECT}}
We define a general left-handed scalar superfield $\Phi$, a right-handed scalar superfield $\Phi^\dagger$, and a vector superfield $V$ in the Wess-Zumino gauge by\footnote{
    $\Phi$, $\Phi^\dagger$ and $V$ are functions of the superspace coordinates $(x,\theta,\bar\theta)$, and their component fields are functions of the spacetime coordinate $x$. These dependencies are omitted in order to avoid cluttered notation.
}
\begin{subequations}
\label{eq:superfields}
\begin{align}
    \Phi &= A + i\theta\sigma^\mu\bar\theta\partial_\mu A - \frac{1}{4}\thetafour\Box A
    + \sqrt{2}\theta\psi - \frac{i}{\sqrt{2}}\theta\theta\partial_\mu\psi\sigma^\mu\bar\theta + \theta\theta F,
    \label{eq:scalar_LH}
    \\
    \Phi^\dagger &= A^* - i\theta\sigma^\mu\bar\theta\partial_\mu A^* - \frac{1}{4}\thetafour\Box A^*
    + \sqrt{2}\bar\theta\bar\psi + \frac{i}{\sqrt{2}}\bar\theta\bar\theta\theta{\sigma^\mu}\partial_\mu\bar\psi + \bar\theta\bar\theta F^*,
    \label{eq:scalar_RH}
    \\
    V &= \theta\sigma^\mu\bar\theta V_\mu
    + \theta\theta\bar\theta\bar\lambda + \bar\theta\bar\theta\theta\lambda + \thetafour D,
    \label{eq:vector_WZ}
\end{align}
\end{subequations}
where $A$ and $F$ are complex scalars, $\psi$ and $\lambda$ are Weyl spinors, $V_\mu$ is a real four-vector, and $D$ is a real scalar.
$A$, $F$, $\psi$, $\lambda$, $V_\mu$ and $D$ are collectively called \emph{component fields}.
$\theta^A$ and $\bar\theta_{\dot A}$ with $A,\dot A = 1,2$ are Grassmann coordinates
with contractions defined by $\theta\theta \equiv \theta^A\theta_A$, $\bar\theta\bar\theta \equiv \bar\theta_{\dot A}\bar\theta^{\dot A}$, and $\theta\sigma^\mu\bar\theta \equiv \theta^A\sigma^\mu_{A \dot A}\bar\theta^{\dot A}$, where $\sigma^\mu=(1,\vec{\sigma})$ and $\vec\sigma$ is a vector whose components are the Pauli matrices.

We define the supersymmetric Lagrangian $\tilde{\mathcal{L}}$, which is a function of superfields, such that it relates to the standard field theory Lagrangian $\mathcal{L}$ by
% $\mathcal{L} = \int d^4\theta \, \tilde{\mathcal{L}}$,
\begin{equation}
    \mathcal{L} = \int d^4\theta \, \tilde{\mathcal{L}},
\end{equation}
where $\int d^4\theta$ acts as a projection operator leaving only terms in $\tilde{\mathcal{L}}$ proportional to $\thetafour$.
A general supersymmetric Lagrangian consisting of scalar superfields $\Phi_i$ and vector superfields $V_j$ is given as
\begin{equation}
    \tilde{\mathcal{L}} = K[\Phi_i^\dagger,\Phi_i,V_j] + \bar\theta\bar\theta W[\Phi_i] + \theta\theta W[\Phi_i^\dagger] + \bar\theta\bar\theta F[V_j],
    \label{eq:general_lagrangian}
\end{equation}
where $K$ is the \emph{K\"ahler potential}, $W$ is the superpotential and $F$ gives the supersymmetric field strength terms.
Note in particular that the K\"ahler potential is the only term that combines vector- and scalar superfields, and thus the only term that can contain slepton--gauge boson interactions.


\subsubsection{The K\"ahler potential \textcolor{red}{FROM SUSY PROJECT}}\label{sec:kahler}
We write the K\"ahler potential for a scalar superfield $\Phi$ transforming in the fundamental representation of some gauge group $G$ as
\begin{equation}
    \tilde{\mathcal L}_\text{K\"ahler} = \Phi^\dagger e^{-2qV^aT^a} \Phi, \label{eq:kahler}
\end{equation}
where $V^a$ are vector superfields needed to keep the K\"ahler potential supergauge invariant, $T_a$ are the generators of $G$, and $q$ is the charge of $\Phi$ under $G$.
Since $\Phi$ transforms under the fundamental representation of $G$ it is generally a multiplet of scalar superfields.
We define the supergauge transformation of a (left-handed) scalar field as
\begin{equation}
    \Phi \to e^{2iq\Lambda^a T^a}\Phi \label{eq:gauge_scalar},
\end{equation}
where $\Lambda^a$ are also (left-handed) scalar superfields.
To keep the K\"ahler potential supergauge invariant we must therefore require that the vector superfield transforms according to
\begin{equation}
    e^{-2qV^aT_a} \to e^{2iq\Lambda^{a\dagger} T^a} e^{-2qV^a T^a} e^{-2iq\Lambda^a T^a}.
\end{equation}
Using the Baker-Campbell-Hausdorff formula \rednote{ref/cite} we can show that this means the vector superfield itself transforms as
\begin{equation}
    V^a \to V^a + i\bigl(\Lambda^a - \Lambda^{a\dagger} \bigr) + qf^{abc} V^b (\Lambda^c + \Lambda^{c\dagger}) + \mathcal{O}(\Lambda^2).
    \label{eq:gauge_vector}
    % V &\to V + i\bigl( \Lambda-\Lambda^\dagger \bigr) - i \bigl[ V,\Lambda+\Lambda^\dagger \bigr] + \cdots, \label{eq:gauge_vector}
\end{equation}


\subsection{Supersymmetry Breaking}


\section{The Minimal Supersymmetric Standard Model}
\subsection{Field Content and Lagrangian Terms}
\subsection{Slepton Phenomenology \textcolor{gray}{(Mass terms, mixing, decay to neutralinos, missing momentum, detectors, current state of parameter space, current state-of-the art corrections?)}}

\subsubsection{Mixing of eigenstates \textcolor{red}{FROM SUSY PROJECT}}
$W^3_\mu$ and $B_\mu$ mix to give the mass eigenstates $Z_\mu$ (the massive Z boson) and $A_\mu$ (the massless photon), while $W^1_\mu$ and $W^2_\mu$ mix to give the charged W-bosons $W^\pm_\mu$. These mixings are given by
\begin{subequations}
    \begin{align}
        Z_\mu &= c_W W^3_\mu - s_W B_\mu, \label{eq:ew_Z}
        \\
        A_\mu &= c_W B_\mu + s_W W^3_\mu, \label{eq:ew_A}
        \\
        W^\pm_\mu &= \frac{1}{\sqrt{2}} \bigl( W^1_\mu \mp iW^2_\mu \bigr) \label{eq:ew_W},
    \end{align}
\end{subequations}
where $c_W \equiv \cos\theta_W$ and $s_W \equiv \sin\theta_W$, and $\theta_W$ is the Weinberg angle.
The electric charge is related to the weak charge $g$ and hypercharge $g'$ by
\begin{equation*}
    e = \frac{gg'}{\sqrt{g^2+g'^2}} = gs_W = g'c_W.
\end{equation*}
Some useful relations that follow directly from these definitions are
\begin{subequations}
    \begin{gather}
        W^{1\mu}W^1_\mu + W^{2\mu}W^2_\mu = 2W^{+\mu}W^-_\mu, \label{eq:W12_squared}
        \\
        Z_\mu = \frac{1}{\sqrt{g^2+g'^2}} \bigl( gW^3_\mu - g'B_\mu \bigr),
        \\
        A_\mu = \frac{1}{\sqrt{g^2+g'^2}} \bigl( g'W^3_\mu + gB_\mu \bigr).
    \end{gather}
\end{subequations}


Similarly, the left- and right-handed charged sleptons mix to give the mass eigenstates $\tilde\ell_{iA}$ (the states that diagonalize the mass matrix), where $A=1,2$. The sneutrinos, having no right-handed counterpart, do not mix.
We introduce the slepton mixing matrices $L^i_{\alpha A}$ defined by
\begin{equation}
    \begin{pmatrix}
        \tilde\ell_{iL} \\ \tilde\ell_{iR}
    \end{pmatrix} = \underbrace{\begin{pmatrix}
        c_{\tilde\ell_i} & s_{\tilde\ell_i}^* \\
        -s_{\tilde\ell_i} & c_{\tilde\ell_i}
    \end{pmatrix}}_{L^i_{\alpha A}} \begin{pmatrix}
        \tilde\ell_{i1} \\ \tilde\ell_{i2}
    \end{pmatrix},
\end{equation}
where $i$ denotes the slepton generation, $\alpha\in\{L,R\}$ the gauge eigenstate, and $A\in\{1,2\}$ the mass eigenstate, and its elements are restricted by $|c_{\tilde\ell_i}|^2+|s_{\tilde\ell_i}|^2=1$.
This definition is similar to the one in~\cite{Richardson:2000nt}, but here we allow for complex values in the off-diagonal elements.
$L^i_{\alpha A}$ is unitary, and it follows that
\begin{align*}
    L^{i*}_{LA}L^i_{LB} + L^{i*}_{RA}L^i_{RB}
    &= \delta_{AB},
\end{align*}
i.e.\@ $\tilde\ell_{iL}^*\tilde\ell_{iL}+\tilde\ell_{iR}^*\tilde\ell_{iR} = \delta_{AB}\tilde\ell_{iA}\tilde\ell_{iB}$.
Following~\cite{Richardson:2000nt} we define the following slepton couplings
\begin{subequations}
\begin{align}
    \zliab &= \frac{1}{2}\bigl( L^{i*}_{LA}L^i_{LB} - 2s_W^2 \delta_{AB} \bigr),
    \label{eq:zliab}
    \\
    \znuiab &= -\frac{1}{2}\delta_{A=1,B=1}.
    \label{eq:znuiab}
\end{align}
\end{subequations}


%%%%%%%%%%%%%%%%%%%%%
%%% Feynman rules %%%
%%%%%%%%%%%%%%%%%%%%%
\subsection{Feynman Rules for Sleptons}\label{sec:feynman_rules}


\subsubsection{General Scalar--Vector Component Interactions from Superfields}
We want to find the interaction between sleptons (the scalar component fields of the scalar superfields) and the electroweak gauge bosons (the vector component fields in the vector superfields), thus for our discussion most components in the superfields are irrelevant.
Considering the general superfields in \cref{eq:superfields} the relevant terms remaining are\footnote{
    Note that we only include terms from $\Phi$ with at most one $\theta$ and one $\bar\theta$ since any other term will give zero when multiplied with the vector superfield.
}
\begin{align*}
    \Phi &\supset A + i\theta\sigma^\mu\bar\theta \partial_\mu A,
    \\
    \Phi^\dagger &\supset A^* - i\theta\sigma^\mu\bar\theta \partial_\mu A^*,
    \\
    V &\supset \theta\sigma^\mu\bar\theta V_\mu.
\end{align*}
As we saw from \cref{eq:general_lagrangian}, any slepton--gauge boson interaction term must come from the K\"ahler potential.
When Taylor expanding the exponential in \cref{eq:kahler} only the terms up to quadratic order survives since, in the Wess-Zumino gauge, we have that $V^3=0$.
Using indices $i,j,k$ and $\ell$ to denote different superfields we see that there are two types of interaction terms that can appear in the K\"ahler potential: three-point interactions, given by
\begin{align}
    V_i\Phi_j^\dagger\Phi_k &\supset \theta\sigma^\mu\bar\theta V_{i\mu}
    \Bigl[ A_j^* - i\theta\sigma^\nu\bar\theta\partial_\nu A_j^* \Bigr]
    \Bigl[ A_k + i\theta\sigma^\rho\bar\theta\partial_\rho A_k \Bigr]
    \nonumber
    \\
    \underset{\thetafour}&{\supset} \frac{i}{2} \thetafour \, V_i^\mu \Bigl( A_j^*\lrpartial{\mu} A_k \Bigr), \label{eq:general_three_point}
\end{align}
with $g \lrpartial{\mu} f \equiv g (\partial_\mu f) - (\partial_\mu g)f$;
and four-point interactions, given by
\begin{align}
    V_iV_j\Phi_k^\dagger\Phi_\ell &\supset (\theta\sigma^\mu\bar\theta)V_{i\mu}
    (\theta\sigma^\nu\bar\theta)V_{j\nu}
    A_k^*A_\ell \nonumber
    \\
    &= \frac{1}{2}\thetafour \, V_i^\mu V_{j\mu} A_k^*A_\ell. \label{eq:general_four_point}
\end{align}
\Cref{eq:general_three_point,eq:general_four_point} will be utilized frequently throughout the following discussion.


\subsubsection{Feynman Rules for Slepton--Gauge Boson Interactions}
Using that the generators of $SU(2)_L$ are $\frac{1}{2}\sigma_a$,
and that $L_i$ and $\bar E_i$ have hypercharges $Y=-1$ and $Y=2$,\footnote{
    In the $Q=T_3+\frac{1}{2}Y$ convention. \rednote{Define earlier in theory (section on SM?)}
} respectively,
it is fairly straightforward to write down the relevant terms in the MSSM K\"ahler potential using \cref{eq:kahler}:
\begin{equation*}
    \tilde{\mathcal{L}} \supset L_i^\dagger e^{-g \sigma_a W^a + g'B} L_i + \bar{E}_i^\dagger e^{-2g'B} \bar{E}_i.
\end{equation*}
Expanding the exponential and ignoring the kinetic terms gives
\begin{equation}
    \tilde{\mathcal{L}} \supset
    \underbrace{-2g'B \bar{E}_i^\dagger \bar{E}_i}_{\equiv \tilde{\mathcal{L}}_{3R}}
    + \underbrace{L_i^\dagger (-g\sigma_a W^a + g'B) L_i}_{\equiv \tilde{\mathcal{L}}_{3L}}
    + \underbrace{2g'^2 BB \bar{E}_i^\dagger \bar{E}_i}_{\equiv \tilde{\mathcal{L}}_{4R}}
    + \underbrace{\frac{1}{2} L_i^\dagger (g\sigma_a W^a - g'B)^2 L_i}_{\equiv \tilde{\mathcal{L}}_{4L}}.
    \label{eq:mssm_kahler}
\end{equation}


\paragraph{Three-Point Interactions}\label{sec:three_point}
The three-point interactions come from $\tilde{\mathcal L}_{3R}$ and $\tilde{\mathcal L}_{3L}$ in the first line of \cref{eq:mssm_kahler}. To find them in terms of component fields we expand $\tilde{\mathcal L}_{3L}$:
\begin{equation}
    \tilde{\mathcal{L}}_{3L} =
    -g (W^1-iW^2)\nu_i^\dagger\ell_i - g(W^1+iW^2)\ell_i^\dagger\nu_i
    + (g'B-gW^3)\nu_i^\dagger \nu_i
    + (g'B+gW^3)\ell_i^\dagger\ell_i.
    \label{eq:lagrange_3L}
\end{equation}
All the interactions in terms of component fields can then be found using \cref{eq:general_three_point} and integrating over the Grassmann coordinates.
Three-point interactions involving two charged sleptons come from $\tilde{\mathcal{L}}_{3R}$ and the last term of $\tilde{\mathcal{L}}_{3L}$ in \cref{eq:lagrange_3L}:
\begin{align*}
    \mathcal{L}_{V\tilde\ell\tilde\ell} &= -ig' B^\mu \Bigl( \tilde{\ell}_{iR} \lrpartial{\mu} \tilde{\ell}_{iR}^* \Bigr)
    +\frac{i}{2} (gW^{3\mu}+g'B^\mu) \Bigl( \tilde{\ell}_{iL}^* \lrpartial{\mu} \tilde{\ell}_{iL} \Bigr)
    \\
    &= ie \biggl( A^\mu - \frac{s_W^2}{s_Wc_W}Z^\mu \biggr) \Bigl( \tilde{\ell}_{iR}^* \lrpartial{\mu} \tilde{\ell}_{iR} \Bigr)
    % \\ &\quad
    + \biggl[ ie A^\mu + \frac{ie}{2s_Wc_W} (1-2s_W^2) Z^\mu \biggr] \Bigl( \tilde{\ell}_{iL}^* \lrpartial{\mu}  \tilde{\ell}_{iL} \Bigr),
\end{align*}
which can be rewritten in terms of the sleptons' mass eigenstates, giving
\begin{align}
    \mathcal{L}_{V\tilde\ell\tilde\ell}
    &= \biggl[ ieA^\mu \delta^{AB} + \frac{ie}{s_Wc_W} \zliab Z^\mu \biggr] \tilde{\ell}_{iA}^* \lrpartial{\mu} \tilde{\ell}_{iB},
    \label{eq:L_Vll}
\end{align}
where $A$ and $B$ label the sleptons' mass eigenstates.
Three-point interactions involving two sneutrinos come from the third term of $\tilde{\mathcal{L}}_{3L}$ in \cref{eq:lagrange_3L}:
\begin{align}
    \mathcal{L}_{V\tilde\nu\tilde\nu} &= \frac{i}{2} (g'B_\mu-gW^3_\mu) \Bigl( \tilde\nu_{iL}^* \lrpartialUpper{\mu} \tilde\nu_{iL} \Bigr)
    \nonumber
    \\
    &= \frac{ie}{s_Wc_W} \znuiab Z_\mu \Bigl( \tilde\nu_{iA}^* \lrpartialUpper{\mu} \tilde\nu_{iB} \Bigr).
    \label{eq:L_Vnunu}
\end{align}
The remaining terms, the first two terms of $\tilde{\mathcal{L}}_{3L}$ in \cref{eq:lagrange_3L}, give the three-point interactions involving both a charged slepton and a sneutrino. The Lagrangian terms are
\begin{align}
    \mathcal{L}_{V\tilde\ell\tilde\nu} &= -\frac{ig}{2} \biggl[
        (W^1_\mu-iW^2_\mu) \Bigl( \tilde\nu_{iL}^* \lrpartialUpper{\mu} \tilde\ell_{iL} \Bigr)
        + (W^1_\mu + iW^2_\mu) \Bigl( \tilde\ell_{iL}^* \lrpartialUpper{\mu} \tilde\nu_{iL} \Bigr)
    \biggr]
    \nonumber
    \\
    &= -\frac{ie}{\sqrt{2}s_W} \biggl[
        W^+_\mu \Bigl( \tilde\nu_{iL}^* \lrpartialUpper{\mu} \tilde\ell_{iL} \Bigr)
        + W^-_\mu \Bigl( \tilde\ell_{iL}^* \lrpartialUpper{\mu} \tilde\nu_{iL} \Bigr)
    \biggr].
    \label{eq:L_Vlnu}
\end{align}

The resulting Feynman rules from the Lagrangian terms in \cref{eq:L_Vll,eq:L_Vlnu,eq:L_Vnunu} are given in \cref{tab:slepton_rules}.
We can see that they coincide perfectly with Refs.~\cite{Richardson:2000nt,Rosiek:1989sce}.

\paragraph{Four-Point Interactions}
The four-point interactions come from $\tilde{\mathcal L}_{4R}$ and $\tilde{\mathcal L}_{4L}$ in the second line of \cref{eq:mssm_kahler}.
Using the anticommutator relation for Pauli matrices:
$\{\sigma_a,\sigma_b\} = 2 \delta_{ab}$,
it is easy to show that
$(\sigma_a W^a)^2 = W_aW^a$.
Using this we can expand $\tilde{\mathcal{L}}_{4L}$ to get\footnote{
    Here we place the indices of $W$ downstairs, reserving the superscripts for exponents.
}
\begin{equation}
    \begin{split}
        \tilde{\mathcal{L}}_{4L} &=
        \frac{1}{2} \Bigl[
            g^2 (W_1^2+W_2^2) + (gW_3+g'B)^2
        \Bigr] \ell_i^\dagger \ell_i
        \\ &\quad
        + \frac{1}{2} \Bigl[
            g^2 (W_1^2+W_2^2) + (gW_3 - g'B)^2
        \Bigr] \nu_i^\dagger \nu_i
        \\ &\quad
        -gg'B \Bigl[
            (W_1-iW_2)\nu_i^\dagger\ell_i + (W_1+iW_2)\ell_i^\dagger\nu_i
        \Bigr].
        \label{eq:lagrange_4L}
    \end{split}
\end{equation}
The interaction terms can then be found by applying \cref{eq:general_four_point} and integrating over the Grassmann coordinates.
Four-point interactions involving two charged sleptons come from $\tilde{\mathcal{L}}_{4R}$ and the first term of $\tilde{\mathcal{L}}_{4L}$ in \cref{eq:lagrange_4L}:
\begin{align*}
    \mathcal{L}_{VV\tilde\ell\tilde\ell}
    &= \frac{1}{4} \Biggl[
        \frac{2e^2}{s_W^2} W^{+\mu}W^-_\mu + 4e^2 A^\mu A_\mu
        + \frac{4e^2}{s_Wc_W}(1-2s_W^2) A^\mu Z_\mu
        \\ &\quad
        + \frac{e^2}{s_W^2c_W^2}(1 - 2s_W^2)^2 Z^\mu Z_\mu
    \Biggr] \tilde\ell_{iL}^*\tilde\ell_{iL}
    + e^2 \biggl( A^\mu A_\mu - \frac{2s_W}{c_W} A^\mu Z_\mu + \frac{s_W^2}{c_W^2} Z^\mu Z_\mu \biggr)
    \tilde\ell_{iR}^*\tilde\ell_{iR},
\end{align*}
which in terms of mass eigenstates is
\begin{equation}
    \begin{split}
        \mathcal{L}_{VV\tilde\ell\tilde\ell}
        &= \tilde\ell_{iA}^*\tilde\ell_{iB} \Biggl[
            \frac{e^2}{2s_W^2} W^{+\mu}W^-_\mu L^{i*}_{LA}L^i_{LB}
            + e^2 \delta^{AB} A^\mu A_\mu
            \\ & \quad
            + \frac{2e^2}{s_Wc_W} \zliab A^\mu Z_\mu
            + \frac{e^2}{c_W^2} \Biggl(
                \frac{1-4s_W^2}{4s_W^2} L_{LA}^{i*}L_{LB}^i + s_W^2 \delta^{AB}
            \Biggr) Z^\mu Z_\mu
        \Biggr].
    \end{split}
    \label{eq:L_VVll}
\end{equation}
The terms involving two sneutrinos come from the second line of $\tilde{\mathcal{L}}_{4L}$ in \cref{eq:lagrange_4L}:
\begin{align}
    \hspace{-0.5em} \mathcal{L}_{VV\tilde\nu\tilde\nu}
    &= \frac{1}{4} \Biggl[
        2\frac{e^2}{s_W^2} W^{+\mu}W^-_\mu
        + \frac{e^2}{s_W^2c_W^2} Z^\mu Z_\mu
    \Biggr] \tilde\nu_i^*\tilde\nu_i,
    \label{eq:L_VVnunu}
\end{align}
and the terms involving both a charged slepton and a sneutrino come from the last term of $\tilde{\mathcal{L}}_{4L}$ in \cref{eq:lagrange_4L}:
\begin{equation}
    \hspace{-0.5em} \mathcal{L}_{VV\tilde\ell\tilde\nu}
    = -\frac{e^2}{\sqrt{2}} \biggl(
        \frac{1}{s_W} A^\mu - \frac{1}{c_W} Z^\mu
    \biggr)
    (W^+_\mu \tilde\nu_{iL}^*\tilde\ell_{iL} + W^-_\mu \tilde\ell_{iL}^*\tilde\nu_{iL}).
    \label{eq:L_VVlnu}
\end{equation}

The resulting Feynman rules from the Lagrangian terms in \cref{eq:L_VVll,eq:L_VVlnu,eq:L_VVnunu} are given in \cref{tab:slepton_rules} (note the symmetry factor of $2$ from couplings with $Z^\mu Z_\mu$ and $A^\mu A_\mu$).
They are exactly the same as the ones found in Ref.~\cite{Rosiek:1989sce}\footnote{
    We are not able to cross-check these rules with Richardson's thesis~\cite{Richardson:2000nt} as he does not include four-point interactions.
}


\begin{table}[ht]
    \centering
    \caption{Feynman rules for slepton--gauge boson interactions. $V^i_\mu$ and $s_i$ denote general vectors and charged scalars.}
    \label{tab:slepton_rules}
    \renewcommand{\arraystretch}{1.8}
    \begin{tabular}{|c|c|c|}
        \hline
        \multicolumn{3}{|c|}{\textbf{Three-point interactions}}
        \\ \hline
        Coupling & Feynman rule & Diagram
        \\ \hline
        $\gamma\tilde\ell_{iA}^*\tilde\ell_{iB}$ & $ie (p+p')^\mu \delta^{AB}$ & \multirow{5}{*}{\inputtikz[baseline=(v)]{rules/general-3point}}
        \\
        $Z\tilde\ell_{iA}^*\tilde\ell_{iB}$ & $\frac{ie}{s_Wc_W} Z_{\tilde\ell_{i}}^{AB} (p+p')^\mu$ &
        \\
        $Z\tilde\nu_{iA}^*\tilde\nu_{iB}$ & $\frac{ie}{s_Wc_W} \znuiab (p+p')^\mu$ &
        \\
        $W^-\tilde\ell_{iA}^*\tilde\nu_{iL}$ & $-\frac{ie}{\sqrt{2} s_W} L^{i*}_{LA} (p+p')^\mu$ &
        \\
        $W^+\tilde\nu_{iL}^*\tilde\ell_{iA}$ & $-\frac{ie}{\sqrt{2} s_W} L^i_{LA} (p+p')^\mu$ &
        \\ \hline\hline
        \multicolumn{3}{|c|}{\textbf{Four-point interactions}}
        \\ \hline
        Coupling & Feynman rule & Diagram
        \\ \hline
        $\gamma\gamma \tilde\ell_{iA}^*\tilde\ell_{iB}$ & $2ie^2\delta^{AB}g^{\mu\nu}$ & \multirow{10}{*}{\inputtikz[baseline=(v)]{rules/general-4point}}
        \\
        $\gamma Z \tilde\ell_{iA}^*\tilde\ell_{iB}$ & $\frac{2ie^2}{s_Wc_W} \zliab g^{\mu\nu}$ &
        \\
        $Z Z \tilde\ell_{iA}^*\tilde\ell_{iB}$ & $\frac{2ie^2}{c_W^2} \Bigl( \frac{1-4s_W^2}{4s_W^2} L_{LA}^{i*}L_{LB}^i + s_W^2 \delta^{AB} \Bigr) g^{\mu\nu}$ &
        \\
        $W^+ W^- \tilde\ell_{iA}^*\tilde\ell_{iB}$ & $\frac{ie^2}{2s_W^2} L^{i*}_{LA} L^i_{LB} g^{\mu\nu}$ &
        \\
        $W^+ W^- \tilde\nu_{iL}^*\tilde\nu_{iL}$ & $\frac{ie^2}{2s_W^2} g^{\mu\nu}$ &
        \\
        $Z Z \tilde\nu_{iL}^*\tilde\nu_{iL}$ & $\frac{ie^2}{2s_W^2c_W^2} g^{\mu\nu}$ &
        \\
        $\gamma W^- \tilde\ell_{iA}^*\tilde\nu_{iL}$ & $-\frac{ie^2}{\sqrt{2}s_W} L^{i*}_{LA} g^{\mu\nu}$ &
        \\
        $\gamma W^+ \tilde\nu_{iL}^*\tilde\ell_{iA}$ & $-\frac{ie^2}{\sqrt{2}s_W} L^i_{LA} g^{\mu\nu}$ &
        \\
        $Z W^- \tilde\ell_{iA}^*\tilde\nu_{iL}$ & $\frac{ie^2}{\sqrt{2}c_W} L^{i*}_{LA} g^{\mu\nu}$ &
        \\
        $Z W^+ \tilde\nu_{iL}^*\tilde\ell_{iA}$ & $\frac{ie^2}{\sqrt{2}c_W} L^i_{LA} g^{\mu\nu}$ &
        \\ \hline
    \end{tabular}
\end{table}

% \end{document}